\documentclass[a4paper,11pt]{article}
\usepackage{lmodern}
\usepackage[T1]{fontenc}
\usepackage[utf8]{inputenc}
\usepackage{amsmath,amssymb,amsthm}
\usepackage{mathtools}
\usepackage{booktabs}
\usepackage{longtable}
\usepackage{array}
\usepackage{hyperref}
\usepackage{graphicx}
\usepackage{float}
\usepackage{geometry}
\geometry{margin=25mm}

\hypersetup{
  colorlinks=true,
  linkcolor=blue,
  urlcolor=blue,
  citecolor=blue
}

\theoremstyle{definition}
\newtheorem{theorem}{Theorem}[section]
\newtheorem{proposition}[theorem]{Proposition}
\newtheorem{lemma}[theorem]{Lemma}
\newtheorem{corollary}[theorem]{Corollary}
\newtheorem{definition}[theorem]{Definition}
\newtheorem{remark}[theorem]{Remark}
\newtheorem{observation}[theorem]{Observation}

\setlength{\parindent}{1em}
\setlength{\parskip}{0.5em}

\providecommand{\tightlist}{\setlength{\itemsep}{0pt}\setlength{\parskip}{0pt}}
\providecommand{\real}[1]{#1}
\begin{document}

\section{\texorpdfstring{Gnomonic Projection from \(S^4\) and the
Schwarzschild-like Form of the Induced Metric on a Four-Dimensional
Subjective Chart (Paper 2: Mathematical Foundations of the
Projection)}{Gnomonic Projection from S\^{}4 and the Schwarzschild-like Form of the Induced Metric on a Four-Dimensional Subjective Chart (Paper 2: Mathematical Foundations of the Projection)}}\label{gnomonic-projection-from-s4-and-the-schwarzschild-like-form-of-the-induced-metric-on-a-four-dimensional-subjective-chart-paper-2-mathematical-foundations-of-the-projection}

\textbf{Author}: Noriaki Kihara (WF System Co., Ltd.) \textbf{ORCID}:
\href{https://orcid.org/0009-0004-6753-4020}{0009-0004-6753-4020}
\textbf{Date}: April 2026 \textbf{DOI}:
\href{https://doi.org/10.5281/zenodo.19839394}{10.5281/zenodo.19839394}

\begin{center}\rule{0.5\linewidth}{0.5pt}\end{center}

\subsection{Abstract}\label{abstract}

We study the gnomonic (central) projection from a four-dimensional
sphere \(S^4(R) \subset \mathbb{R}^5\) of radius \(R\) onto a
four-dimensional tangent hyperplane, and the metric induced on the
hyperplane by pull-back of the standard metric on \(S^4(R)\). We derive
the explicit form of the induced metric, compute its Christoffel
symbols, Riemann tensor, and Einstein tensor, and show that it satisfies
the vacuum Einstein equations with a positive cosmological constant
\(\Lambda = 3/R^2\) in four dimensions. Restricted to a fixed-\(R\)
slice and analytically continued to Lorentzian signature, the
construction yields a metric of de Sitter-like form, and a comparison
with the four-plus-one-dimensional Schwarzschild--Tangherlini solution
reveals a structural similarity in the horizon behaviour. The paper is
mathematical; physical interpretation is deferred.

\begin{center}\rule{0.5\linewidth}{0.5pt}\end{center}

\subsection{§1. Introduction}\label{introduction}

The gnomonic projection from a sphere onto a tangent plane has the
elementary property that great circles on the sphere map to straight
lines on the plane. In two dimensions this is a classical observation in
cartography. In higher dimensions the analogous construction maps great
\((n-1)\)-circles on \(S^n\) to \((n-1)\)-flats on the tangent
hyperplane, and the resulting metric on the hyperplane --- obtained by
pulling back the standard metric on the sphere along the projection ---
is a non-trivial Riemannian metric whose curvature reflects the radius
\(R\) of the sphere.

This paper studies the four-dimensional case in detail: gnomonic
projection from \(S^4(R) \subset \mathbb{R}^5\) to a four-dimensional
tangent hyperplane \(\Pi_R \subset \mathbb{R}^5\). We derive the induced
metric in explicit coordinates, compute its curvature invariants, and
identify a Schwarzschild-like form of the metric in spherical
coordinates. The principal results are:

\begin{enumerate}
\def\labelenumi{\arabic{enumi}.}
\tightlist
\item
  The induced metric satisfies \(G_{\mu\nu} + \Lambda\, g_{\mu\nu} = 0\)
  with \(\Lambda = 3/R^2\) in four dimensions, a positive cosmological
  constant.
\item
  After Lorentzian continuation, the metric is the de Sitter metric in
  Beltrami coordinates.
\item
  In spherical coordinates with the radial coordinate identified
  appropriately, the metric reduces to a Schwarzschild--de Sitter form,
  which in the limit of small mass parameter coincides with pure de
  Sitter and in the limit of vanishing cosmological constant
  (\(R \to \infty\)) coincides with pure Schwarzschild.
\item
  The natural higher-dimensional comparison object is the
  four-plus-one-dimensional Schwarzschild--Tangherlini solution, with
  which the present construction shares the \(r^{-2}\) falloff of the
  gravitational potential.
\end{enumerate}

Sections are organized as follows. Section 2 sets up the gnomonic
projection and the tangent-hyperplane coordinates. Section 3 derives the
induced metric and computes the Christoffel symbols. Section 4 computes
the Riemann and Einstein tensors. Section 5 analyzes the metric in
spherical coordinates and identifies the Schwarzschild-like form.
Section 6 compares with the Tangherlini solution. Section 7 concludes.

The paper is purely mathematical. The interpretation of the construction
in terms of physical black hole thermodynamics is the subject of
separate work.

\begin{center}\rule{0.5\linewidth}{0.5pt}\end{center}

\subsection{\texorpdfstring{§2. The Construction: Gnomonic Projection
from
\(S^4\)}{§2. The Construction: Gnomonic Projection from S\^{}4}}\label{the-construction-gnomonic-projection-from-s4}

\subsubsection{§2.1 Setup}\label{setup}

Let \(\mathbb{R}^5\) be the ambient space with coordinates
\(Y = (Y^0, Y^1, Y^2, Y^3, Y^4)\) and standard Euclidean metric
\(\delta_{AB}\, dY^A \otimes dY^B\). The four-dimensional sphere of
radius \(R\) is
\[S^4(R) = \{Y \in \mathbb{R}^5 : (Y^0)^2 + (Y^1)^2 + (Y^2)^2 + (Y^3)^2 + (Y^4)^2 = R^2\}.\]

We choose the \(Y^0\)-axis as the projection axis. The tangent
hyperplane at the ``north pole'' \((R, 0, 0, 0, 0)\) is
\[\Pi_R = \{Y \in \mathbb{R}^5 : Y^0 = R\},\] parametrized by
coordinates \((x^1, x^2, x^3, x^4)\) via \(Y^\mu = x^\mu\) for
\(\mu = 1, 2, 3, 4\).

\subsubsection{§2.2 The projection map}\label{the-projection-map}

The gnomonic projection \(\Phi : \Pi_R \to S^4(R)_+\) (the ``northern
hemisphere'' \(Y^0 > 0\)) is defined by the rule that the line through
the origin and a point \(\Phi(x) \in S^4(R)\) also passes through
\(x \in \Pi_R\). Explicitly:
\[\Phi(x) = \frac{R}{\ell(x)} (R, x^1, x^2, x^3, x^4), \qquad \ell(x) := \sqrt{R^2 + |x|^2}, \quad |x|^2 := \sum_{\mu=1}^{4} (x^\mu)^2.\]

One verifies directly that \(\sum_{A=0}^{4} \Phi(x)^A \Phi(x)^A = R^2\),
so \(\Phi(x) \in S^4(R)\).

\subsubsection{§2.3 Inverse and domain of
regularity}\label{inverse-and-domain-of-regularity}

The inverse map \(\Phi^{-1} : S^4(R)_+ \to \Pi_R\) is
\[\Phi^{-1}(Y)^\mu = \frac{R Y^\mu}{Y^0} \qquad (\mu = 1, 2, 3, 4).\]

The map is well-defined and smooth on the open hemisphere \(Y^0 > 0\).
The equator \(Y^0 = 0\) is mapped to ``infinity'' in the sense that
\(\Phi^{-1}(Y)\) diverges as \(Y^0 \to 0^+\). The southern hemisphere
\(Y^0 < 0\) is not in the image of \(\Phi\); for the southern points one
can use a separate projection through the south pole, but we will not
need to do so here.

\subsubsection{§2.4 Geodesic image}\label{geodesic-image}

Great circles (geodesics) on \(S^4(R)\) are the intersections of
\(S^4(R)\) with two-dimensional linear subspaces of \(\mathbb{R}^5\)
through the origin. Since these subspaces are stable under the rescaling
\(Y \mapsto \lambda Y\), their images under \(\Phi^{-1}\) are the
intersections of these subspaces with \(\Pi_R\), which are straight
lines (or affine flats) in \(\Pi_R\). This is the defining property of
the gnomonic projection: geodesics on the sphere correspond to straight
lines on the tangent hyperplane.

\begin{center}\rule{0.5\linewidth}{0.5pt}\end{center}

\subsection{§3. The Induced Metric}\label{the-induced-metric}

\subsubsection{§3.1 Pull-back computation}\label{pull-back-computation}

The standard metric on \(S^4(R)\) is the restriction of the ambient
Euclidean metric. Pulling back along \(\Phi\), we obtain a metric \(g\)
on \(\Pi_R\):
\[g_{\mu\nu}(x) = \delta_{AB}\, \frac{\partial \Phi^A}{\partial x^\mu}\, \frac{\partial \Phi^B}{\partial x^\nu}.\]

A direct computation gives
\[\frac{\partial \Phi^0}{\partial x^\mu} = -\frac{R^2 x_\mu}{\ell^3}, \qquad \frac{\partial \Phi^\nu}{\partial x^\mu} = \frac{R}{\ell} \delta^\nu_\mu - \frac{R x^\nu x_\mu}{\ell^3}.\]

Substituting and simplifying:
\[\boxed{\;g_{\mu\nu}(x) = \frac{R^2}{\ell^2}\left(\delta_{\mu\nu} - \frac{x_\mu x_\nu}{\ell^2}\right).\;}\]

\subsubsection{§3.2 Inverse metric}\label{inverse-metric}

Direct calculation (or matrix inversion using the Sherman-Morrison
formula):
\[g^{\mu\nu}(x) = \frac{\ell^2}{R^2}\left(\delta^{\mu\nu} + \frac{x^\mu x^\nu}{R^2}\right).\]

The relation \(g_{\mu\nu} g^{\nu\rho} = \delta_\mu^\rho\) is verified by
computation.

\subsubsection{§3.3 Determinant}\label{determinant}

Using
\(\det(\delta_{\mu\nu} - x_\mu x_\nu / \ell^2) = 1 - |x|^2/\ell^2 = R^2/\ell^2\):
\[\det g = \left(\frac{R^2}{\ell^2}\right)^4 \cdot \frac{R^2}{\ell^2} = \frac{R^{10}}{\ell^{10}}.\]

The volume element is \(\sqrt{\det g}\, d^4 x = (R/\ell)^5 d^4 x\).

\subsubsection{§3.4 Christoffel symbols}\label{christoffel-symbols}

The Christoffel symbols
\(\Gamma^\rho_{\mu\nu} = \frac{1}{2} g^{\rho\sigma}(\partial_\mu g_{\nu\sigma} + \partial_\nu g_{\mu\sigma} - \partial_\sigma g_{\mu\nu})\)
can be computed in closed form:
\[\Gamma^\rho_{\mu\nu} = -\frac{1}{\ell^2}\left(x^\rho \delta_{\mu\nu} + (\text{symmetric corrections})\right).\]

The complete expression is given in Appendix A. \textbf{(後で検討: full
expression with the cross-term contributions.)}

\begin{center}\rule{0.5\linewidth}{0.5pt}\end{center}

\subsection{§4. Riemann and Einstein
Tensors}\label{riemann-and-einstein-tensors}

\subsubsection{§4.1 Riemann tensor}\label{riemann-tensor}

Computing the Riemann tensor from the Christoffel symbols:
\[R^\rho_{\sigma\mu\nu} = \frac{1}{R^2}(g_{\sigma\mu} \delta^\rho_\nu - g_{\sigma\nu} \delta^\rho_\mu).\]

This is the Riemann tensor of a maximally symmetric space ---
specifically, of a constant-positive-curvature space with sectional
curvature \(1/R^2\).

\subsubsection{§4.2 Ricci tensor and scalar
curvature}\label{ricci-tensor-and-scalar-curvature}

Contracting:
\[R_{\sigma\nu} = R^\mu_{\sigma\mu\nu} = \frac{3}{R^2} g_{\sigma\nu}, \qquad R = g^{\mu\nu} R_{\mu\nu} = \frac{12}{R^2}.\]

\subsubsection{§4.3 Einstein tensor and the cosmological
constant}\label{einstein-tensor-and-the-cosmological-constant}

The Einstein tensor in \(n = 4\) dimensions:
\[G_{\mu\nu} := R_{\mu\nu} - \frac{1}{2} R\, g_{\mu\nu} = \frac{3}{R^2} g_{\mu\nu} - \frac{6}{R^2} g_{\mu\nu} = -\frac{3}{R^2} g_{\mu\nu}.\]

Therefore \(G_{\mu\nu} + \Lambda g_{\mu\nu} = 0\) with
\[\boxed{\;\Lambda = \frac{3}{R^2}.\;}\]

The induced metric satisfies the four-dimensional vacuum Einstein
equations with a positive cosmological constant whose value is
determined by the radius of the sphere.

\subsubsection{§4.4 Lorentzian
continuation}\label{lorentzian-continuation}

Continuing one of the coordinates to Lorentzian signature --- say,
replacing \(x^4 \to i\tau\) --- converts the Riemannian metric on
\(\Pi_R\) to a Lorentzian metric. Direct substitution shows that the
resulting metric is the \textbf{de Sitter metric in Beltrami
coordinates}:
\[ds^2 = \frac{R^2}{\ell_L^2}\left(-d\tau^2 + |dx|^2 + (\text{cross terms})\right), \qquad \ell_L^2 := R^2 - \tau^2 + |x|^2.\]

The cosmological constant in this Lorentzian setting is
\(\Lambda = 3/R^2 > 0\), and the metric is the four-dimensional de
Sitter spacetime in unconventional coordinates.

\begin{center}\rule{0.5\linewidth}{0.5pt}\end{center}

\subsection{§5. Spherical Coordinates and the Schwarzschild-like
Form}\label{spherical-coordinates-and-the-schwarzschild-like-form}

\subsubsection{§5.1 Spherical
decomposition}\label{spherical-decomposition}

In Lorentzian signature, decompose the spatial part using spherical
coordinates \((r, \theta, \phi)\) with \(r = |\vec{x}|\):
\[ds^2 = -A(\tau, r)\, d\tau^2 + B(\tau, r)\, dr^2 + r^2\, d\Omega_2^2 + (\text{cross terms}),\]
where the explicit forms of \(A, B\) follow from §4.4.

\subsubsection{§5.2 Static slicing}\label{static-slicing}

For a particular slicing (constant-\(\tau\) surfaces appropriately
chosen --- \textbf{後で検討}: precise gauge choice), the metric reduces
to
\[ds^2 = -f(r)\, d\tau^2 + f(r)^{-1} dr^2 + r^2\, d\Omega_2^2, \qquad f(r) = 1 - \frac{r^2}{R^2}.\]

This is the de Sitter metric in static coordinates, with cosmological
horizon at \(r = R\).

\subsubsection{§5.3 Adding a mass
parameter}\label{adding-a-mass-parameter}

A natural extension introduces a mass parameter \(M\) via
\[f(r) = 1 - \frac{2GM}{r} - \frac{r^2}{R^2}.\]

This is the \textbf{Schwarzschild--de Sitter} metric, with both an event
horizon (near \(r = 2GM\)) and a cosmological horizon (near \(r = R\)).

In the limit \(M \to 0\), the metric reduces to pure de Sitter
(cosmological horizon at \(r = R\)). In the limit \(R \to \infty\), the
metric reduces to pure Schwarzschild (event horizon at \(r = 2GM\)).

The presence of a mass-like term in the four-dimensional metric is a
candidate physical source for the central-projection construction;
whether this term has a microscopic origin in the discrete-cube packing
of the companion paper is left open.

\begin{center}\rule{0.5\linewidth}{0.5pt}\end{center}

\subsection{§6. Comparison with the Four-Plus-One-Dimensional
Schwarzschild--Tangherlini
Solution}\label{comparison-with-the-four-plus-one-dimensional-schwarzschildtangherlini-solution}

\subsubsection{§6.1 The Tangherlini
metric}\label{the-tangherlini-metric}

The four-plus-one-dimensional Schwarzschild--Tangherlini metric is
\[ds^2 = -f_T(r)\, dt^2 + f_T(r)^{-1} dr^2 + r^2\, d\Omega_3^2, \qquad f_T(r) = 1 - \frac{8GM}{3\pi r^2}.\]

The key features are: - The horizon at \(r_h = (8GM/(3\pi))^{1/2}\). -
The radial fall-off \(1/r^2\) in the gravitational potential,
characteristic of \(D = 5\). - The Hawking temperature
\(T_H = 1/(2\pi r_h)\). - The Bekenstein--Hawking entropy
\(S = 2\pi^2 r_h^3/(4 G_5)\).

\subsubsection{§6.2 Structural similarity with the central-projection
metric}\label{structural-similarity-with-the-central-projection-metric}

The central-projection construction of §5 produces a metric of the form
\(1 - r^2/R^2\) at the cosmological horizon, and (in the
Schwarzschild--de Sitter extension) of the form \(1 - 2GM/r - r^2/R^2\).
The 4+1 Tangherlini metric, by contrast, has \(1 - 8GM/(3\pi r^2)\).

The two metrics are not equivalent. They share the qualitative feature
of having a horizon, and in particular the curvature radius \(R\) in the
central-projection construction plays a role analogous to the horizon
radius \(r_h\) in the Tangherlini construction. The precise
correspondence is mediated by the identification
\(R \leftrightarrow r_h\), which in turn is supported by the dynamical
analysis of the companion paper on energy-radius scaling.

\subsubsection{§6.3 Domain of validity}\label{domain-of-validity}

The structural similarity holds at the level of the metric form and the
existence of a horizon. It does not extend to a full identity of the
metrics. The companion paper (Paper 4) develops the quantitative
comparison via the volume deficit \(\Delta(R)\) and the
Bekenstein--Hawking entropy.

\begin{center}\rule{0.5\linewidth}{0.5pt}\end{center}

\subsection{§7. Conclusion}\label{conclusion}

We have derived the gnomonic projection metric on a four-dimensional
tangent hyperplane to \(S^4(R)\), computed its Riemann and Einstein
tensors, and identified its structural form. The principal results are:

\begin{enumerate}
\def\labelenumi{\arabic{enumi}.}
\tightlist
\item
  The induced metric satisfies \(G_{\mu\nu} + (3/R^2) g_{\mu\nu} = 0\).
\item
  Lorentzian continuation gives the four-dimensional de Sitter metric.
\item
  In static spherical coordinates with a mass parameter, the metric is
  the Schwarzschild--de Sitter form.
\item
  The construction shares structural features with, but is not identical
  to, the four-plus-one-dimensional Schwarzschild--Tangherlini metric.
\end{enumerate}

The construction provides the geometric foundation for Hypothesis 1 of
the research program: that black holes in the projected four-dimensional
theory may be the images of higher-dimensional structures. The specific
form of the induced metric and its relation to Tangherlini are the
content of this paper; the physical interpretation, including the
dynamical determination of \(R\) from the matter content of the
four-dimensional theory, is the subject of Paper 5.

\begin{center}\rule{0.5\linewidth}{0.5pt}\end{center}

\subsection{Appendix A: Christoffel Symbols
(deferred)}\label{appendix-a-christoffel-symbols-deferred}

\textbf{後で検討}: Full closed-form expression of
\(\Gamma^\rho_{\mu\nu}\) for the central-projection metric.

\subsection{References}\label{references}

\begin{enumerate}
\def\labelenumi{\arabic{enumi}.}
\tightlist
\item
  K. Schwarzschild, ``Über das Gravitationsfeld eines Massenpunktes nach
  der Einsteinschen Theorie,'' \emph{Sitzungsber. Preuss. Akad. Wiss.},
  189--196 (1916).
\item
  F. Kottler, ``Über die physikalischen Grundlagen der Einsteinschen
  Gravitationstheorie,'' \emph{Annalen der Physik} \textbf{56}, 401
  (1918).
\item
  F. R. Tangherlini, ``Schwarzschild Field in \(n\) Dimensions and the
  Dimensionality of Space Problem,'' \emph{Nuovo Cimento} \textbf{27},
  636 (1963).
\item
  J. D. Bekenstein, ``Black Holes and Entropy,'' \emph{Phys. Rev.~D}
  \textbf{7}, 2333 (1973).
\item
  S. W. Hawking, ``Particle Creation by Black Holes,'' \emph{Comm. Math.
  Phys.} \textbf{43}, 199 (1975).
\item
  R. C. Myers, M. J. Perry, ``Black Holes in Higher Dimensional
  Spacetimes,'' \emph{Annals Phys.} \textbf{172}, 304 (1986).
\end{enumerate}

\end{document}
